\section{总体方案与微架构设计}

\subsection{系统总体架构}

\texttt{YH\_rv\_cpu} 采用 ``CPU 核 + SoC 封装 + FPGA 原型顶层'' 的层次化组织方式。
其中 CPU 内核负责指令执行与异常控制，SoC 层负责连接同步指令存储、数据存储和最小 MMIO 外设，
FPGA 顶层负责将工程映射到 Nexys A7-100T 的原型实现路径。

\begin{table}[H]
    \centering
    \caption{核心模块分工}
    \begin{tabular}{p{0.28\textwidth} p{0.58\textwidth}}
        \toprule
        模块/文件 & 作用 \\
        \midrule
        \texttt{rtl/YH\_rv\_cpu.v} & CPU 顶层，组织五级流水、前递、redirect、CSR 与 trap \\
        \texttt{rtl/YH\_rv\_cpu\_hazard\_unit.v} & load-use hazard 检测与前递控制 \\
        \texttt{rtl/YH\_rv\_cpu\_soc.v} & SoC 顶层，连接 ROM、RAM、UART、timer、DONE \\
        \texttt{rtl/YH\_rv\_sync\_imem\_rom.v} & 指令 ROM，服务同步取指路径 \\
        \texttt{rtl/YH\_rv\_dmem\_ram.v} & 数据 RAM，服务 load/store 路径 \\
        \texttt{fpga/vivado/src/YH\_rv\_cpu\_fpga\_top.v} & FPGA 顶层与综合参数入口 \\
        \bottomrule
    \end{tabular}
\end{table}

\FigurePlaceholder{系统总体架构图占位：CPU/SoC/FPGA 顶层关系}

\subsection{CPU 五级流水组织}

当前 CPU 主线采用经典五级流水：

\begin{enumerate}
    \item IF：同步取指、PC 更新、redirect 生效；
    \item ID：译码、寄存器读取、立即数生成、控制信号准备；
    \item EX：ALU、分支跳转判断、CSR 访问、trap/mret/redirect 决策；
    \item MEM：load/store 访存与 MMIO 访问；
    \item WB：将 ALU、访存或 CSR 结果写回寄存器堆。
\end{enumerate}

这种结构的优点是边界清晰、验证成本可控，且足以支撑当前比赛阶段的 CoreMark 和基础软件运行。
与单纯追求复杂架构不同，本项目更强调可回归性和文档口径稳定性，因此当前并未在主线中引入
更高风险的多发射、乱序或 cache 复杂机制。

\subsection{数据相关处理}

为了维持五级流水的吞吐与正确性，项目实现了明确的数据相关处理策略：

\begin{itemize}
    \item EX/MEM、MEM/WB 两级前递，降低通用 ALU 相关带来的停顿；
    \item \texttt{load-use hazard} 检测，当加载结果尚未可用且后继指令立即消费时，对译码级施加停顿；
    \item 当前正式保留的停顿策略为 \texttt{stall\_decode = load\_use\_hazard}，
    这是经过 fresh CoreMark 与回归验证确认后的保留优化。
\end{itemize}

这种策略的直接特点是：对可前递场景尽量放行，对不可前递且风险明确的 load-use 场景做最小必要停顿，
在稳定性和性能之间取得平衡。

\subsection{控制相关与 redirect 机制}

当前控制流处理主要依赖 EX 级决策：

\begin{itemize}
    \item branch/jump/jalr 在 EX 级形成 redirect；
    \item \texttt{trap}、\texttt{mret} 与普通跳转共享控制流重定向框架；
    \item 当 \texttt{ex\_fetch\_redirect\_valid} 有效时，IF 级更新 PC，清空错误路径气泡。
\end{itemize}

2026-04-09 的 split profile 显示，当前 CoreMark 的 redirect 开销以 taken branch 为主，
其中 \texttt{bne} 和 \texttt{beq} 占据了主要比例；这说明控制流路径的代价已经比 request-side
微调更值得关注。不过项目在初赛提交版中仍保持保守策略，不把尚未证明收益的新控制流改动写进主线。

\FigurePlaceholder{流水线数据通路与 redirect/trap 路径图占位}

\subsection{CSR、异常与中断处理}

项目当前实现的控制状态寄存器与异常处理，重点覆盖比赛阶段确有需求的机器态功能：

\begin{itemize}
    \item CSR 指令：\texttt{csrrw/csrrs/csrrc} 及其立即数变体；
    \item 机器态关键 CSR：\texttt{mstatus}、\texttt{mie}、\texttt{mtvec}、
    \texttt{mscratch}、\texttt{mepc}、\texttt{mcause}、\texttt{mip}；
    \item 同步异常：非法指令、CSR 非法访问、load misaligned 等；
    \item 控制流相关：\texttt{ecall}、\texttt{ebreak}、\texttt{mret}；
    \item 机器定时器中断：通过 \texttt{timer} 外设触发并进入 trap 流程。
\end{itemize}

赛题答疑已经明确 ``优先完成核心指令即可''，因此当前实现边界是工程上合理的：
既保证关键软件与测试可运行，也避免为初赛阶段引入无验证支撑的高复杂度特权体系。

\subsection{SoC 存储与 MMIO 组织}

SoC 采用 ``ROM + RAM + MMIO'' 的最小闭环结构。其中同步 ROM 负责取指，RAM 负责数据访问，
MMIO 外设负责提供对外可观测能力和计时能力。当前已确认的地址组织如下：

\begin{table}[H]
    \centering
    \caption{SoC 关键地址映射}
    \begin{tabular}{p{0.28\textwidth} p{0.24\textwidth} p{0.32\textwidth}}
        \toprule
        资源 & 地址/基址 & 作用 \\
        \midrule
        ROM & \texttt{0x0000\_0000} & 程序存储基址 \\
        RAM & \texttt{0x0000\_4000} & 数据存储基址 \\
        UART TX & \texttt{0x1000\_0000} & 串口发送寄存器 \\
        DONE & \texttt{0x1000\_0004} & 结束标志寄存器 \\
        TIMER VALUE LO/HI & \texttt{0x1000\_0008}/\texttt{0x1000\_000C} & 当前计时值 \\
        TIMER CMP LO/HI & \texttt{0x1000\_0010}/\texttt{0x1000\_0014} & 比较寄存器 \\
        TIMER CTRL & \texttt{0x1000\_0018} & 定时器控制与中断使能 \\
        \bottomrule
    \end{tabular}
\end{table}

这种组织方式具有三个优势：

\begin{itemize}
    \item 对仿真、CoreMark 与 \texttt{riscv-tests} 都足够直接；
    \item 对 FPGA 原型移植成本低，易于保持路径一致；
    \item 可以把 UART、timer 和 DONE 作为最小观测接口，便于后续实板 bring-up。
\end{itemize}

\subsection{当前保留优化与已关闭方向}

当前正式保留的优化只有三项：

\begin{enumerate}
    \item \texttt{stall\_decode = load\_use\_hazard}
    \item FPGA 默认 \texttt{IMEM\_OUTPUT\_REG = 0}
    \item FPGA 默认 \texttt{DMEM\_OUTPUT\_REG = 0}
\end{enumerate}

与之对应，\texttt{request-cursor}、\texttt{pipe-hit}、redirect 同拍取指、
\texttt{FQ-01 \textasciitilde{} FQ-06A} 等方向均已完成实验或 quick-screen，并被证明
没有形成可保留收益，或者带来过高风险。因此，本说明书只保留经验证能稳定复现的结果，
不将未保留试验混入正式设计结论。

\clearpage
